\documentclass[12pt,a4paper]{article}
\usepackage[margin=25mm, headheight=15pt, headsep=10pt]{geometry}
\usepackage{fontspec}
\usepackage[table]{xcolor} % Note: table option is required for rowcolors
\usepackage{titlesec}
\usepackage{tocloft}
\usepackage{fancyhdr}
\usepackage{booktabs}
\usepackage{tcolorbox}
\tcbuselibrary{skins, breakable}
\usepackage[colorlinks=true, linkcolor=emerald, urlcolor=emerald, bookmarks=true]{hyperref}

\definecolor{emerald}{RGB}{0,102,51}
\definecolor{darkred}{RGB}{139,0,0}
\definecolor{lightgray}{RGB}{245,245,245}

\usepackage[nil]{babel}
\babelprovide[import, main]{bengali}
\babelprovide[import]{arabic}
\babelfont{rm}[Renderer=HarfBuzz,Scale=1.0]{Noto Serif Bengali}
\babelfont{sf}[Renderer=HarfBuzz]{Noto Sans Bengali}
\babelfont{tt}[Renderer=HarfBuzz]{Noto Sans Bengali}
\babelfont[arabic]{rm}[Renderer=HarfBuzz,Scale=1.1]{Noto Naskh Arabic}

% Section styling
\titleformat{\section}{\Large\bfseries\color{emerald}}{\thesection.}{0.5em}{}
\titleformat{\subsection}{\large\bfseries\color{emerald!85!black}}{\thesubsection.}{0.5em}{}
\titleformat{\subsubsection}{\normalsize\bfseries\color{emerald!70!black}}{\thesubsubsection.}{0.5em}{}
\titlespacing*{\section}{0pt}{18pt}{8pt}

% Fancy Header/Footer setup
\pagestyle{fancy}
\fancyhf{} % clear all header and footer fields
\fancyhead[L]{\color{emerald}\small\textbf{\nouppercase{\leftmark}}}
\fancyfoot[L]{\scriptsize\color{black!50}{\lat{AI}}-সংকলিত, আলেম-যাচাইকৃত নয়}
\fancyfoot[C]{\color{emerald!70!black}\thepage}
\renewcommand{\headrulewidth}{0.4pt}
\renewcommand{\headrule}{\hbox to\headwidth{\color{emerald}\leaders\hrule height \headrulewidth\hfill}}
\renewcommand{\footrulewidth}{0pt}

% Custom boxes
% 1. Ayat/Hadith Box
\newtcolorbox{ayatbox}[1][]{
    enhanced, breakable, 
    colback=emerald!5, 
    colframe=emerald, 
    boxrule=0.5pt, 
    arc=3pt, 
    left=10pt, right=10pt, top=10pt, bottom=10pt,
    #1
}

% 2. Quote/Translation Box
\newtcolorbox{quotebox}[1][]{
    enhanced, breakable,
    frame hidden,
    colback=lightgray,
    borderline west={3pt}{0pt}{emerald},
    left=15pt, right=10pt, top=10pt, bottom=10pt,
    sharp corners,
    #1
}

% 3. AI-generated notice box
\newtcolorbox{warnbox}[1][]{
    enhanced, breakable,
    colback=red!4,
    colframe=darkred,
    boxrule=0.8pt,
    arc=3pt,
    left=10pt, right=10pt, top=10pt, bottom=10pt,
    #1
}

% Commands
\newcommand{\arab}[1]{\foreignlanguage{arabic}{#1}}
\newcommand{\kib}[1]{\textcolor{emerald}{\textbf{#1}}}

% Latin (English) fallback: Noto Serif Bengali has NO Latin glyphs
\newfontfamily\latfont[Renderer=HarfBuzz]{Noto Serif}
\newcommand{\lat}[1]{{\latfont #1}}

\setlength{\parskip}{5pt}
\setlength{\parindent}{0pt}

% Fix for missing bullet in Noto Serif Bengali
\renewcommand{\labelitemi}{$\bullet$}



\begin{document}
\begin{center}{\Large\bfseries\color{emerald} \lat{13}-নামাজের-পর-যিকর-দোয়া}\end{center}
\vspace{1em}
\section*{১৩ — নামাজের পরের যিকর ও দোয়া}

\begin{warnbox}
 \textbf{গুরুত্বপূর্ণ নোটিশ:} এই কন্টেন্টটি \lat{AI} দিয়ে প্রস্তুত। কেবল সহীহ/হাসান হাদিস রাখা হয়েছে।
\end{warnbox}

\medskip\hrule\medskip

\subsection*{১. সূত্র কার্ড}

\begin{table}[h]\centering\small
\begin{tabular}{ll}
বিষয় & বিবরণ \\
\hline
\textbf{বিষয়} & নামাজের পর ইস্তিগফার, তসবিহ ফাতিমা, দোয়া \\
\textbf{সূত্র} & সহীহ মুসলিম ৫৯১, ৫৯৬ \\
\textbf{মর্যাদা} & \textbf{সুন্নাতে মুআক্কাদাহ} \\
\hline
\end{tabular}
\end{table}

\medskip\hrule\medskip

\subsection*{২. তসবিহ ফাতিমা — মূল আরবি}

\textbf{সুবহানাল্লাহ} (৩৩ বার), \textbf{আলহামদুলিল্লাহ} (৩৩ বার), \textbf{আল্লাহু আকবার} (৩৩ বার) — মোট ৯৯ বার। এরপর: \textbf{লা ইলাহা ইল্লাল্লাহু ওয়াহদাহু লা শারীকা লাহ, লাহুল মুলকু ওয়া লাহুল হামদু ওয়া হুয়া আলা কুল্লি শাই'িন কাদির} (১ বার) — মোট ১০০।

\medskip\hrule\medskip

\subsection*{৩. শব্দ-শব্দ অর্থ}

\begin{table}[h]\centering\small
\begin{tabular}{lll}
আরবি & বাংলা & \lat{English} \\
\hline
\textbf{\arab{سُبْحَانَ} \arab{اللَّهِ}} & আল্লাহ পবিত্র & \textit{\lat{Glory} \lat{be} \lat{to} \lat{Allah}} \\
\textbf{\arab{الْحَمْدُ} \arab{لِلَّهِ}} & সকল প্রশংসা আল্লাহর & \textit{\lat{Praise} \lat{be} \lat{to} \lat{Allah}} \\
\textbf{\arab{اللَّهُ} \arab{أَكْبَرُ}} & আল্লাহ সবচেয়ে বড় & \textit{\lat{Allah} \lat{is} \lat{the} \lat{Greatest}} \\
\textbf{\arab{لَا} \arab{إِلَٰهَ} \arab{إِلَّا} \arab{اللَّهُ}} & আল্লাহ ছাড়া ইলাহ নেই & \textit{\lat{There} \lat{is} \lat{no} \lat{god} \lat{but} \lat{Allah}} \\
\textbf{\arab{وَحْدَهُ} \arab{لَا} \arab{شَرِيكَ} \arab{لَهُ}} & একক, শরীক নেই & \textit{\lat{Alone}, \lat{no} \lat{partner}} \\
\textbf{\arab{لَهُ} \arab{الْمُلْكُ}} & তাঁরই রাজত্ব & \textit{\lat{To} \lat{Him} \lat{dominion}} \\
\textbf{\arab{وَلَهُ} \arab{الْحَمْدُ}} & তাঁরই প্রশংসা & \textit{\lat{and} \lat{praise}} \\
\textbf{\arab{وَهُوَ} \arab{عَلَىٰ} \arab{كُلِّ} \arab{شَيْءٍ} \arab{قَدِيرٌ}} & সব কিছুর ওপর ক্ষমতাবান & \textit{\lat{He} \lat{has} \lat{power} \lat{over} \lat{all}} \\
\hline
\end{tabular}
\end{table}

\medskip\hrule\medskip

\subsection*{৪. পূর্ণ বাংলা অনুবাদ}

\textbf{\lat{“}আল্লাহ পবিত্র\lat{”} (৩৩ বার), \lat{“}সকল প্রশংসা আল্লাহর\lat{”} (৩৩ বার), \lat{“}আল্লাহ সবচেয়ে বড়\lat{”} (৩৩ বার)। এরপর: \lat{“}আল্লাহ ছাড়া কোনো ইলাহ নেই — একক, কোনো শরীক নেই। তাঁরই রাজত্ব, তাঁরই প্রশংসা, সব কিছুর ওপর তিনি ক্ষমতাবান।\lat{”} (১ বার)}

\medskip\hrule\medskip

\subsection*{৫. ইস্তিগফার (নামাজের পর)}

\textbf{\lat{“}আস্তাগফিরুল্লাহা\lat{”}} (৩ বার), এরপর: \textbf{\lat{“}আস্তাগফিরুল্লাহা ওয়া আতুবু ইলাইহি\lat{”}} (১ বার)

\textbf{অর্থ:} \lat{“}আমি আল্লাহর কাছে ক্ষমা চাই এবং তাঁর দিকে তাওবা করি।\lat{”}

\medskip\hrule\medskip

\subsection*{৬. সংক্ষিপ্ত ব্যাখ্যা}

নামাজের পর সোজা বসে ডান হাতের আঙুল দিয়ে গোনুন। ডান দিকে ইস্তিগফার, বাম দিকে দরুদ। এরপর যা দোয়া চান তা চাইতে পারেন — এটি \textbf{\lat{accepted} \lat{time} (কবুল হওয়ার সময়)}।

\medskip\hrule\medskip

\subsection*{৭. সংশ্লিষ্ট হাদিস}

\begin{table}[h]\centering\small
\begin{tabular}{ll}
হাদিস & গ্রন্থ \\
\hline
\textbf{\lat{“}নামাজের পর 'সুবহানাল্লাহ' ৩৩ বার, 'আলহামদুলিল্লাহ' ৩৩ বার, 'আল্লাহু আকবার' ৩৩ বার বলো।\lat{”}} & সহীহ মুসলিম ৫৯১ \\
\textbf{\lat{“}নামাজের পর 'আস্তাগফিরুল্লাহ' ৩ বার বলো।\lat{”}} & সহীহ মুসলিম ৫৯১ \\
\hline
\end{tabular}
\end{table}

\medskip\hrule\medskip

\textit{শেষ আপডেট: আগস্ট ২০২৬}


\end{document}
